\section{Introduction}
\label{sec:introduction}

Fantasy football games translate on-pitch match events into a numeric
score per player per match. The official FIFA Fantasy game for the
2026 World Cup follows the template familiar from Fantasy Premier
League: managers select a 15-player squad under a budget cap, field a
starting eleven in a valid formation, and nominate a captain whose
points are doubled; non-starters are replaced through automatic
substitutions, and per-stage rules govern transfers and the number of
players allowed from any one country. The nominal objective is to
maximise total fantasy points over the tournament. In practice most
entries also compete in a private league, which layers a ranking
objective on top of raw scoring, and the two objectives diverge in
non-trivial ways: an entry trailing the field benefits from variance,
accepting differential picks with lower expected value, while an
entry near the lead benefits from template picks that lock in the
field's expected output.

Predicting fantasy points for an international tournament is
structurally harder than the club-football problem addressed by most
prior work \citep{groos2025openfpl, mishra2025fpl}. A World Cup
compresses 104 matches into five weeks, against 380 matches over ten
months in an English Premier League (EPL) season. National teams
assemble for roughly ten days a year, so international form is
observed sparsely. The opponent distribution resembles no club
league: Argentina plays Cape Verde, not a domestic third-tier
equivalent. The group-stage cadence also induces rotation behaviour
that club play does not exhibit: a coach whose team has already
qualified rests starters on the final group matchday. A model trained
purely on club football therefore transfers imperfectly. The features
that drive club fantasy points (opponent strength, home advantage,
position-specific scoring rates) generalise; the features specific to
tournament play (manager rotation policy, top-scorer-race incentives,
knockout-stage attacking emphasis) do not exist in the club training
set at all.

The transfer problem is compounded by a hard resource constraint:
before the opening match there exist zero rows of FIFA Fantasy World
Cup 2026 training data. Our supervised learner is therefore trained
on three seasons of community-archived EPL fantasy scoring
\citep{vaastav}, \EplTrainingRows{}
player-gameweek rows in total, and international signal is injected
at inference time rather than learned. Concretely, we maintain a live
Elo rating \citep{elo1978rating, hvattum2010using, gilch2018elo}
computed over a public archive of every international match since
1872 \citep{martj42}, and the rating difference between a player's
country and its opponent enters every backend as an
opponent-strength feature.

This paper presents the system we built and then operated live for
the full tournament. Expected points are produced by an ensemble of
three backends of increasing statistical ambition: a hand-tuned
heuristic over price, position, matchup, and venue; a structural
Poisson goals model in the tradition of \citet{maher1982modelling}
and \citet{dixon1997modelling}, with positional goal and assist
shares; and a LightGBM regressor \citep{ke2017lightgbm} trained on
the EPL corpus. Squad selection, stage-aware transfer planning under
the official penalty of $-3$ points per extra transfer, and lineup
and captain choice are formulated as a mixed-integer linear
programme, solved with PuLP and the CBC solver \citep{pulp, cbc}.
The constraint set encodes the official rules stage by stage: a
\$100M budget in the group stage rising to \$105M from the round of
32, a per-country cap that relaxes from 3 to 8 as teams are
eliminated, and a rolling free-transfer allowance.

Evaluation is deliberately two-fold. Before the tournament, each
backend's per-position RMSE was measured on a held-out EPL 2024--25
end-of-season split, fixing an honest per-position choice of backend
before any World Cup data existed. During the tournament, we froze
predictions ahead of every round and logged every decision (transfer,
captain, lineup) with its rationale at decision time and its realised
outcome afterwards, including the cases in which a human domain prior
overrode the model, and whether the override was vindicated.
\POSTFINAL{headline result after the final}

This work makes the following contributions:

\begin{enumerate}
\item \textbf{System.} An end-to-end, reproducible pipeline for
  World Cup fantasy management: data collectors with schema
  validation, a per-(player, round) feature table, a three-backend
  expected-points ensemble with a live international Elo signal, and
  a MILP optimiser covering squad selection, stage-aware transfer
  planning with hit penalties, and lineup and captain choice.
\item \textbf{Cross-domain transfer.} An empirical account of
  transferring club-trained fantasy models to international
  tournament play under a zero-shot constraint, including
  per-position held-out validation used to select backends and the
  documented failure modes of the transfer, such as blindness to
  elite players without recent EPL history.
\item \textbf{In-tournament validation protocol.} A pre-registered
  style of live evaluation: predictions frozen before each round, a
  complete decision log with ex-ante rationale and ex-post outcome,
  and an explicit accounting of model recommendations versus human
  domain priors.
\item \textbf{Negative results.} We report what did not work with
  the same care as what did, including blending prediction-market
  prices into expected points \citep{benter1994computer,
  wolfers2004prediction}, a team Elo difference training feature,
  and goalkeeper defensive-form adjustments, each rejected by
  measured validation error.
\item \textbf{Open source and data.} The full codebase, frozen data
  snapshots, and generated evaluation artefacts are released
  publicly, so every table and figure in this paper can be
  regenerated from the repository
  (\POSTFINAL{repository URL once public}).
\end{enumerate}

The remainder of the paper is organised as follows.
Section~\ref{sec:related_work} reviews related work in football
modelling and fantasy points prediction. Section~\ref{sec:data}
documents the five data sources and the harmonisation required to
join them. Section~\ref{sec:methodology} develops the three predictor
backends, the Elo signal, and the MILP formulation.
Section~\ref{sec:implementation} describes the implementation and
reproducibility measures. Section~\ref{sec:validation} reports
held-out validation on the EPL split. Section~\ref{sec:live-results}
logs the live tournament round by round. Section~\ref{sec:analysis}
analyses model versus human decisions, and
Section~\ref{sec:negative-results} collects negative results.
Section~\ref{sec:lessons} distils lessons learned,
Section~\ref{sec:future_work} outlines future work, and
Section~\ref{sec:conclusion} concludes.
